\section{Evaluación ISO/IEC 25010:2023}
\label{sec:iso25010}

Esta sección evalúa el sistema según cinco características del modelo
de calidad ISO/IEC 25010:2023 \cite{iso25010:2023}, con métricas medidas de forma directa
sobre el sistema en ejecución, sin proyecciones ni valores estimados.

\subsection{Compatibilidad}

La suite de pruebas de extremo a extremo con Playwright se ejecutó
sobre tres motores de navegador: Chromium, Firefox y WebKit. La
elección de WebKit responde a una limitación práctica del entorno de
desarrollo, que no cuenta con hardware macOS disponible para probar
Safari nativamente; WebKit es el motor de renderizado subyacente de
Safari y constituye la alternativa más cercana disponible para
\emph{continuous integration} en Linux. De las 15 pruebas ejecutadas
por navegador, 9 pasaron de forma consistente en los tres motores,
mostrando el mismo patrón de éxito o falla independientemente del
navegador; las 6 fallas restantes se atribuyeron a interferencia de
red local durante la ejecución y no a incompatibilidades reales entre
motores de renderizado, ya que ninguna de ellas mostró un patrón
distinto entre navegadores. En el ámbito móvil, la compatibilidad se
garantiza fijando \texttt{minSdk = 26} (Android 8.0), cubriendo la
inmensa mayoría de dispositivos activos según las cifras públicas de
distribución de versiones de Android.

\begin{table}[htbp]
\centering
\caption{Resultados de Compatibilidad (ISO/IEC 25010)}
\label{tab:iso-compatibilidad}
\begin{tabular}{lll}
\toprule
Métrica & Medido & Objetivo \\
\midrule
Pruebas E2E consistentes entre navegadores & 9/15 & Suite completa \\
Compatibilidad móvil (Android) & minSdk 26+ & API 26+ \\
\bottomrule
\end{tabular}
\end{table}

\subsection{Seguridad}

% [VILLAMARIN] Completar con base en docs/iso25010/seguridad-ivan.md
% y docs/iso25010/matriz-seguridad-endpoints.csv. Tabla: porcentaje
% de endpoints con verificacion JWT, ausencia de vulnerabilidades del
% top 10 de OWASP, umbral objetivo 100% de endpoints protegidos salvo
% los explicitamente publicos (/api/v1/internal, /actuator/health).

\subsection{Mantenibilidad}
% [PENDIENTE -- ISAIAS]

\subsection{Fiabilidad}

La fiabilidad del módulo de reservas se sustenta en cuatro mecanismos
verificados automáticamente. Primero, la creación y la aprobación guardan
el resultado asociado a una \texttt{Idempotency-Key}; repetir una operación
equivalente devuelve el resultado previo sin duplicar solicitudes,
reservas, historial ni eventos. El uso de una clave con otro actor,
solicitud o carga útil produce un conflicto explícito.

Segundo, la aprobación serializa el acceso al intervalo de agenda mediante
una escritura conflictiva determinista compatible con CockroachDB. La
prueba concurrente con dos solicitudes distintas para el mismo laboratorio
y franja demuestra que se persiste como máximo una reserva. Tercero, la
cancelación de una solicitud aprobada actualiza en una sola transacción la
solicitud y la reserva programada. Finalmente, las pruebas con
Testcontainers validan el esquema y el comportamiento transaccional sobre
una base real, no únicamente mediante repositorios simulados.

La topología E3 configura tres nodos y factor de replicación tres. La
tolerancia efectiva ante la caída controlada de un nodo pertenece al
protocolo experimental del clúster; sus resultados deben citarse desde la
evidencia de ejecución y no inferirse solamente de la configuración.

\begin{table}[htbp]
\centering
\caption{Evidencia de fiabilidad del módulo de reservas}
\label{tab:iso-fiabilidad-reservas}
\begin{tabular}{lll}
\toprule
Propiedad & Evidencia & Resultado \\
\midrule
Idempotencia de creación & Replay con misma clave y carga & Sin duplicado \\
Idempotencia de aprobación & Replay de la misma solicitud & Una reserva \\
Concurrencia & Dos aprobaciones sobre la misma franja & Máximo una reserva \\
Consistencia de cancelación & Solicitud aprobada y reserva programada & Ambas canceladas \\
Regresión automatizada & Suite de Reservas & 110/110 \\
\bottomrule
\end{tabular}
\end{table}

\subsection{Eficiencia de desempeño}

La evaluación de eficiencia está preparada para ejecutarse siempre a
través del \emph{gateway} y no contra el puerto interno de Reservas. El
protocolo contempla lecturas de disponibilidad, solicitudes, reservas y
agenda, además de escrituras controladas para evitar contaminar la base con
operaciones artificiales masivas.

Para cada corrida se deben conservar el total de peticiones, rendimiento en
peticiones por segundo, tasa de fallos, media, mediana o p50, p95, p99,
mínimo y máximo. Las repeticiones previstas permiten calcular media,
desviación estándar e intervalo de confianza del 95\,\%, y generar
diagramas de caja sin selección manual de resultados. Los errores
\texttt{4xx} esperados por autorización o conflicto se separan de los
\texttt{5xx}, que representan fallos internos.

No se fija un SLA arbitrario: los percentiles y el \emph{throughput} deben
reportarse como valores observados. Al no existir en esta versión del
documento una campaña final completa y reproducible asociada a sus CSV, no
se consignan cifras de rendimiento. Esta distinción evita presentar como
evidencia definitiva una configuración, una ejecución aislada o datos de
otra versión del sistema.
